
\documentclass[letter, 11pt]{article}

    %%%%%%%%%%%%%%%%%%%%%%%%%%%%%%
    \oddsidemargin  0.0in
    \evensidemargin 0.0in
    \textwidth      6.8in
    \headheight     0.0in
    \topmargin      0.0in
    \textheight=9.0in
    \parindent=0in
    %%%%%%%%%%%%%%%%%%%%%%%%%%%%%%


\renewcommand{\familydefault}{\sfdefault} % Police 

\usepackage{amsthm,amsfonts,amssymb,amsmath, xcolor, enumitem}
\usepackage{graphicx}   %To include graphics
\usepackage[small,compact]{titlesec}

\usepackage[utf8]{inputenc}

\usepackage[french]{babel}

\usepackage{fancyhdr}
\pagestyle{fancy}

\newcommand{\var}{\mathbb{V}\text{ar}} %Symboles programmés pour simplifier le texte

\newcommand{\E}{\operatorname{E}}
\newcommand{\Var}{\operatorname{Var}}
\newcommand{\bY}{\mathbf{Y}}
\newcommand{\bX}{\mathbf{X}}
\newcommand{\bbeta}{\boldsymbol{\beta}}
\newcommand{\bepsilon}{\boldsymbol{\varepsilon}}
\newcommand{\bH}{\mathbf{H}}
\newcommand{\trans}{^{\mathsf T}}

\usepackage{version} % Pour faire un texte troué pour les étudiants
%\includeversion{etudiant}\excludeversion{professeur} %version etudiant
\includeversion{professeur}\excludeversion{etudiants} %version prof
\newcommand{\etudiant}[1]{\begin{etudiants}#1 \end{etudiants}}
\newcommand{\prof}[1]{\begin{professeur}\color{blue} {\bf Solution :} #1 \end{professeur}}


\fancyhead[L]{STT-2100 \\ Régression}
\fancyhead[R]{A.-A. Vallée \\ Université Laval}

\newcommand{\N}{\mathcal{N}}
\newcommand{\bo}[1]{\boldsymbol#1}
\begin{document}

\vspace*{0.01 cm}
\begin{center}
{\Large Exercices - Chapitre 5}
\end{center}




\begin{enumerate}[leftmargin=*, wide, label=\bf Exercice \arabic*:]

%%%%%%%%%%%%
% Question 1
\item Pour prédire le prix de vente $Y$, en milliers de dollars, d'une maison, on
ajuste le modèle
\[
  \widehat{Y}=-85+0{,}32x_1+18{,}5x_2-7{,}2x_3,
\]
où $x_1$ est la superficie habitable en mètres carrés, $x_2$ le nombre de
salles de bain et $x_3$ la distance du centre-ville en kilomètres.

\begin{enumerate}[label=\alph*)]
  \item Interprétez chacun des trois coefficients de pente dans le contexte.
  \item Interprétez l'ordonnée à l'origine. Cette interprétation est-elle
        utile dans le contexte?
  \item Comparez deux maisons identiques, sauf que la seconde possède une
        salle de bain supplémentaire. Quelle différence de prix le modèle
        prévoit-il?
\end{enumerate}


\prof{  \begin{enumerate}[label=\alph*)]
  \item 
  \begin{itemize}
    \item En maintenant les autres variables constantes, une augmentation de 1~m$^2$ de la superficie habitable est associée
          à une augmentation de  320~\$, du
          prix de vente moyen prédit;
    \item En maintenant les autres variables constantes, une salle de bain supplémentaire est associée à une augmentation de
           18~500~\$, du prix moyen prédit;
    \item En maintenant les autres variables constantes, une augmentation de 1~km de la distance du centre-ville est associée
          à une diminution de 7~200~\$, du prix
          moyen prédit.
  \end{itemize}

  \item L'ordonnée à l'origine $-85$ correspond au prix moyen prédit, en
        milliers de dollars, pour une maison de superficie nulle, sans salle
        de bain et située à 0~km du centre-ville. Ce profil n'est pas réaliste;
        l'interprétation contextuelle de l'ordonnée à l'origine est donc farfelue ici.

  \item Pour des valeurs identiques de $x_1$ et $x_3$, une augmentation d'une
        unité de $x_2$ produit
        \[
          \Delta\widehat{Y}=18{,}5
        \]
        milliers de dollars. Le prix prédit de la seconde maison est donc
        supérieur de 18~500~\$.

\end{enumerate}

}


%%%%%
% Question 2

\item 
Une étude relie la note finale $Y$ d'un étudiant au nombre d'heures d'étude
hebdomadaires $x_1$ et à son résultat à un test diagnostique $x_2$ :
\[
  \E(Y\mid x_1,x_2)=35+1{,}8x_1+0{,}45x_2.
\]

\begin{enumerate}[label=\alph*)]
  \item Interprétez le coefficient de $x_1$.
  \item Pourquoi serait-il incorrect d'affirmer simplement que deux étudiants
        dont le nombre d'heures d'étude diffère de cinq heures auront des notes
        moyennes différant de neuf points?
  \item Donnez deux profils d'étudiants pour lesquels le modèle prévoit la
        même note moyenne.
\end{enumerate}

\prof{
\begin{enumerate}[label=\alph*)]
  \item À résultat identique au test diagnostique, une heure d'étude
        hebdomadaire supplémentaire est associée à une augmentation de
        1,8 point de la note finale moyenne.

  \item Une différence de cinq heures d'étude entraîne une différence prédite
        de $5(1{,}8)=9$ points seulement si les deux étudiants ont le même
        résultat au test diagnostique. En général,
        \[
          \Delta\E(Y)=1{,}8\Delta x_1+0{,}45\Delta x_2.
        \]
        Une différence entre les résultats diagnostiques modifierait donc la
        comparaison. Pour que l'énoncé soit correct, il faudrait préciser que les deux résultats au test diagnostique sont identiques.

  \item Pour obtenir la même note moyenne, il faut que
        \[
          1{,}8\Delta x_1+0{,}45\Delta x_2=0,
          \qquad\text{soit}\qquad
          \Delta x_2=-4\Delta x_1.
        \]
        Par exemple, les profils $(x_1,x_2)=(5,80)$ et $(10,60)$ donnent tous
        les deux
        \[
          \widehat{Y}=35+1{,}8(5)+0{,}45(80)
                     =35+1{,}8(10)+0{,}45(60)=80.
        \]
\end{enumerate}
}


%%%%
% Question 3

\item On observe la réponse $Y$ ainsi que deux variables explicatives $x_1$ et $x_2$
pour quatre unités :
\[
\begin{array}{c|rrrr}
  i   & 1 & 2 & 3 & 4\\ \hline
  Y   & 5 & 8 & 9 & 14\\
  x_1 & 1 & 2 & 3 & 4\\
  x_2 & 0 & 1 & 1 & 2
\end{array}
\]

\begin{enumerate}[label=\alph*)]
  \item Écrivez explicitement $\bY$, $\bX$ et $\bbeta$.
  \item Calculez $\bX\trans\bX$ et $\bX\trans\bY$.
  \item Vérifiez que $\bX\trans\bX$ est inversible.
  \item En utilisant 
        \[
        (\bX\trans\bX)^{-1}=
        \begin{pmatrix}
          11/4&-2&5/2\\
          -2&2&-3\\
          5/2&-3&5
        \end{pmatrix},
        \] calculez $\widehat{\bbeta}=(\bX\trans\bX)^{-1}\bX\trans\bY$ et donnez l'expression du modèle de prédiction estimé.
  \item Calculez les valeurs ajustées et les résidus.
  \item Vérifiez numériquement que $\bX\trans\widehat{\bepsilon}=\mathbf{0}$.
\end{enumerate}

\prof{
\begin{enumerate}[label=\alph*)]
  \item Puisque le modèle comprend une ordonnée à l'origine,
        \[
        \bY=\begin{pmatrix}5\\8\\9\\14\end{pmatrix},\qquad
        \bX=\begin{pmatrix}
          1&1&0\\
          1&2&1\\
          1&3&1\\
          1&4&2
        \end{pmatrix},\qquad
        \bbeta=\begin{pmatrix}\beta_0\\\beta_1\\\beta_2\end{pmatrix}.
        \]

  \item Les produits matriciels sont
        \[
        \bX\trans\bX=
        \begin{pmatrix}
          4&10&4\\
          10&30&13\\
          4&13&6
        \end{pmatrix},
        \qquad
        \bX\trans\bY=
        \begin{pmatrix}36\\104\\45\end{pmatrix}.
        \]

  \item On calcule
        \[
          \det(\bX\trans\bX)=4\neq0.
        \]
        La matrice est donc inversible. 

  \item 
        \[
        \widehat{\bbeta}
        =(\bX\trans\bX)^{-1}\bX\trans\bY
        =\begin{pmatrix}3{,}5\\1\\3\end{pmatrix}.
        \]
        Le modèle ajusté est donc
        \[
          \widehat{Y}=3{,}5+x_1+3x_2.
        \]

  \item Les valeurs ajustées et les résidus sont
        \[
        \widehat{\bY}=\bX\widehat{\bbeta}
        =\begin{pmatrix}4{,}5\\8{,}5\\9{,}5\\13{,}5\end{pmatrix},
        \qquad
        \widehat{\bepsilon}=\bY-\widehat{\bY}
        =\begin{pmatrix}0{,}5\\-0{,}5\\-0{,}5\\0{,}5\end{pmatrix}.
        \]

  \item Finalement,
        \[
        \bX\trans\widehat{\bepsilon}
        =\begin{pmatrix}
          1&1&1&1\\
          1&2&3&4\\
          0&1&1&2
        \end{pmatrix}
        \begin{pmatrix}0{,}5\\-0{,}5\\-0{,}5\\0{,}5\end{pmatrix}
        =\begin{pmatrix}0\\0\\0\end{pmatrix}.
        \]
        Les résidus sont donc orthogonaux à chacune des colonnes de $\bX$.
\end{enumerate}
}

%%%%%
% Question 4

\item On note $\widehat{\bY}=\bX\widehat{\bbeta}$ et
$\widehat{\bepsilon}=\bY-\widehat{\bY}$.

\begin{enumerate}[label=\alph*)]
  \item Montrez que
        $\bX\trans\widehat{\bepsilon}=\mathbf{0}$.
    \item Montrez que $\widehat{\bY}\trans\widehat{\bepsilon}=0$.
  \item Interprétez géométriquement ces résultats.
\end{enumerate}

\prof{
\begin{enumerate}[label=\alph*)]
  \item On a,
        \[
        \bX\trans\widehat{\bepsilon}
        =\bX\trans(\bY-\bX\widehat{\bbeta})
        =\bX\trans\bY-\bX\trans\bX\widehat{\bbeta}
        =\bX\trans\bY-\bX\trans\bX(\bX\trans\bX)^{-1}\bX^\top\bY
        =\mathbf{0}.
        \]

   \item Puisque $\widehat{\bY}=\bX\widehat{\bbeta}$,
        \[
        \widehat{\bY}\trans\widehat{\bepsilon}
        =(\bX\widehat{\bbeta})\trans\widehat{\bepsilon}
        =\widehat{\bbeta}\trans\bX\trans\widehat{\bepsilon}
        =\mathbf{0}.
        \]

  \item Le vecteur $\widehat{\bY}$ appartient à l'espace engendré par les
        colonnes de $\bX$, tandis que $\widehat{\bepsilon}$ appartient à son
        complément orthogonal. Ainsi, $\widehat{\bY}$ est la projection
        orthogonale de $\bY$ sur l'espace colonne de $\bX$ et est orthogonal à $\widehat{\bepsilon}$.
\end{enumerate}
}

%%%%% 
% Question 5

\item Supposons que $\bX$ est de plein rang colonne et posons
\[
  \bH=\bX(\bX\trans\bX)^{-1}\bX\trans.
\]

\begin{enumerate}[label=\alph*)]
  \item Montrez que $\bX\widehat{\bbeta}=\bH\bY$.
  \item Montrez que $\bH$ est symétrique.
  \item Montrez que $\bH$ est idempotente, c'est-à-dire que
        $\bH^2=\bH$.
  \item Montrez que $\mathbf{I}_n-\bH$ est également symétrique et idempotente.
  \item Quel est le rôle géométrique de $\mathbf{I}_n-\bH$?
\end{enumerate}

\prof{\begin{enumerate}[label=\alph*)]
  \item En remplaçant l'estimateur des moindres carrés,
        \[
        \widehat{\bY}=\bX\widehat{\bbeta}
        =\bX(\bX\trans\bX)^{-1}\bX\trans\bY
        =\bH\bY.
        \]

  \item On utilise $(\mathbf{A}\mathbf{B})\trans
        =\mathbf{B}\trans\mathbf{A}\trans$ et le fait que
        $\bX\trans\bX$ ainsi que son inverse sont symétriques :
        \begin{align*}
        \bH\trans
        &=\left[\bX(\bX\trans\bX)^{-1}\bX\trans\right]\trans\\
        &=\bX\left[(\bX\trans\bX)^{-1}\right]\trans\bX\trans\\
        &=\bX(\bX\trans\bX)^{-1}\bX\trans=\bH.
        \end{align*}

  \item On calcule
        \begin{align*}
        \bH^2
        &=\bX(\bX\trans\bX)^{-1}\bX\trans
          \bX(\bX\trans\bX)^{-1}\bX\trans\\
        &=\bX(\bX\trans\bX)^{-1}
          (\bX\trans\bX)(\bX\trans\bX)^{-1}\bX\trans\\
        &=\bX(\bX\trans\bX)^{-1}\bX\trans=\bH.
        \end{align*}

  \item La symétrie suit de
        \[
          (\mathbf{I}_n-\bH)\trans
          =\mathbf{I}_n-\bH\trans=\mathbf{I}_n-\bH.
        \]
        Pour l'idempotence,
        \[
          (\mathbf{I}_n-\bH)^2
          =\mathbf{I}_n-2\bH+\bH^2
          =\mathbf{I}_n-\bH.
        \]

  \item La matrice $\mathbf{I}_n-\bH$ projette $\bY$ sur le complément
        orthogonal de l'espace colonne de $\bX$. En particulier,
        \[
          \widehat{\bepsilon}=(\mathbf{I}_n-\bH)\bY.
        \]
\end{enumerate}
}


\end{enumerate}

\end{document}

